\chapter{บทนำสู่เกษตรอัจฉริยะและสถาปัตยกรรม AIoT สำหรับไม้ผลเมืองร้อน}
\label{chap:ch01}

\section*{แผนบริหารการสอนประจำบทที่ 1}
\addcontentsline{toc}{section}{แผนบริหารการสอนประจำบทที่ 1}

\noindent\textbf{หัวข้อเนื้อหาประจำบท}
\begin{enumerate}
    \item วิวัฒนาการของการเกษตรและแนวคิดเกษตรแม่นยำสูง
    \item สถาปัตยกรรม 4 ชั้นของระบบ AIoT ทางการเกษตร
    \item ความท้าทายเฉพาะทางในการจัดการสวนทุเรียนภาคตะวันออก
\end{enumerate}

\noindent\textbf{วัตถุประสงค์เชิงพฤติกรรม}
\begin{enumerate}
    \item อธิบายหลักการ ทฤษฎี และสถาปัตยกรรมเทคโนโลยีประจำบทที่ 1 ได้อย่างถูกต้อง
    \item ประยุกต์ใช้เครื่องมือ อุปกรณ์เซนเซอร์ และระบบปัญญาประดิษฐ์เพื่อการจัดการสวนทุเรียนได้อย่างมีประสิทธิภาพ
    \item วิเคราะห์ สรุปผลข้อมูล และแก้ไขปัญหาเชิงฟิสิกส์เกษตรและคุณภาพดินได้อย่างเป็นระบบ
\end{enumerate}

\noindent\textbf{การวัดและประเมินผล}
\begin{table}[H]
\centering\small
\begin{tabularx}{\textwidth}{p{0.55\textwidth} p{0.25\textwidth} c}
\toprule
\textbf{เกณฑ์การประเมินผล} & \textbf{เครื่องมือวัดผล} & \textbf{สัดส่วนคะแนน} \\
\midrule
1. ทักษะการปฏิบัติงานและการใช้อุปกรณ์ & แบบประเมินทักษะภาคสนาม & 40\% \\
2. รายงานข้อมูลและการวิเคราะห์ผล & การตรวจสมุดบันทึกและดาต้าชีท & 30\% \\
3. แบบฝึกหัดและคำถามท้ายบท & แบบทดสอบท้ายบทเรียน & 20\% \\
4. การมีส่วนร่วมและความรับผิดชอบ & การเข้าเรียนและการตรงต่อเวลา & 10\% \\
\midrule
\textbf{รวมทั้งสิ้น} & & \textbf{100\%} \\
\bottomrule
\end{tabularx}
\end{table}
\newpage

% ==========================================================
% เนื้อหาบทเรียนประจำบทที่ 1
% ==========================================================

\section{วิวัฒนาการของการเกษตรและแนวคิดเกษตรแม่นยำสูง}
\index{วิวัฒนาการของการเกษตรและแนวคิดเกษตรแม่นยำสูง}

การศึกษาและทำความเข้าใจเกี่ยวกับวิวัฒนาการของการเกษตรและแนวคิดเกษตรแม่นยำสูง (Evolution of Agriculture and Precision Farming) มีความสำคัญอย่างยิ่งในระบบเกษตรอัจฉริยะ โดยมีรายละเอียดสาระสำคัญดังนี้

การเกษตรกรรมของมนุษยชาติได้พัฒนาผ่าน 4 ยุคสมัยหลัก เริ่มจากการทำเกษตรแบบยังชีพ (Agriculture 1.0) ในยุคโบราณ ต่อเนื่องสู่ยุคการใช้เครื่องจักรกลหนักและสารเคมี (Agriculture 2.0) ในช่วงปฏิวัติเขียว และยุคเกษตรกรรมอัตโนมัติ (Agriculture 3.0) ที่เริ่มนำระบบอิเล็กทรอนิกส์และคอมพิวเตอร์มาช่วยควบคุม

ในยุคปัจจุบัน การเกษตรได้ก้าวเข้าสู่ยุคเกษตรกรรม 4.0 หรือ เกษตรอัจฉริยะ (Smart Agriculture) ซึ่งขับเคลื่อนด้วยข้อมูลขนาดใหญ่ (Big Data), การผสานเทคโนโลยีอินเทอร์เน็ตของสรรพสิ่ง (Internet of Things; IoT) เข้ากับปัญญาประดิษฐ์ (Artificial Intelligence; AI) จนเกิดเป็นแนวคิด AIoT (Artificial Intelligence of Things)

เกษตรแม่นยำสูง (Precision Agriculture) มีเป้าหมายเพื่อการจัดการปัจจัยการผลิต ได้แก่ น้ำ ปุ๋ย แสงแดด และสารชีวภัณฑ์ ให้สอดคล้องกับความต้องการของพืชอย่างจำเพาะเจาะจงในแต่ละจุดของแปลงปลูก (Site-specific management) ช่วยลดการสูญเสียทรัพยากร เพิ่มผลตอบแทนต่อไร่ และรักษาดุลยภาพสิ่งแวดล้อม

\section{สถาปัตยกรรม 4 ชั้นของระบบ AIoT ทางการเกษตร}
\index{สถาปัตยกรรม 4 ชั้นของระบบ AIoT ทางการเกษตร}

การศึกษาและทำความเข้าใจเกี่ยวกับสถาปัตยกรรม 4 ชั้นของระบบ AIoT ทางการเกษตร (Four-Layer AIoT Architecture in Agriculture) มีความสำคัญอย่างยิ่งในระบบเกษตรอัจฉริยะ โดยมีรายละเอียดสาระสำคัญดังนี้

ระบบ AIoT สำหรับสวนไม้ผลเมืองร้อนได้รับการออกแบบโครงสร้างการทำงานแบบเป็นลำดับชั้น 4 ระดับ (Four-Layer Architecture):

1. ชั้นการรับรู้และการตรวจวัด (Perception Layer): ประกอบด้วยชุดเซนเซอร์ตรวจวัดสภาพแวดล้อมทางฟิสิกส์และเคมี เช่น เซนเซอร์วัดอุณหภูมิและความชื้นในอากาศ เซนเซอร์วัดรังสีแสงอาทิตย์ และโพรบวัดความชื้นรวมถึงค่าการนำไฟฟ้าในดิน
2. ชั้นเครือข่ายการส่งผ่านข้อมูล (Network Layer): ทำหน้าที่ส่งสัญญาณข้อมูลจากเซนเซอร์ภาคสนามเข้าสู่ระบบเครือข่าย โดยใช้โพรโทคอลไร้สายกำลังต่ำระยะไกล เช่น LoRaWAN, NB-IoT และ Wi-Fi Long Range
3. ชั้นการประมวลผลเอดจ์และคลาวด์ (Edge and Cloud Processing Layer): นำข้อมูลดิบมาผ่านการกรอง ประมวลผลล่วงหน้า และวิเคราะห์เชิงลึกด้วยอัลกอริทึมปัญญาประดิษฐ์ เช่น แบบจำลองคำนวณการคายระเหยน้ำ และโครงข่ายประสาทเทียมตรวจวินิจฉัยสุขภาพดิน
4. ชั้นการประยุกต์ใช้งานและควบคุม (Application Layer): หน้าจอแดชบอร์ดแสดงผลข้อมูลแบบเรียลไทม์ ระบบสั่งการเปิด-ปิดวาล์วน้ำอัตโนมัติ และการแจ้งเตือนสภาวะวิกฤตผ่านสมาร์ทโฟนของเกษตรกร

\begin{techbox}[ข้อกำหนดทางเทคนิคสถาปัตยกรรม AIoT]
การเชื่อมต่อระบบเซนเซอร์ไร้สายต้องรักษาความมั่นคงปลอดภัยของข้อมูล มีการเข้ารหัสระดับ AES-128 และบริหารจัดการพลังงานให้โหนดเซนเซอร์ทำงานได้อย่างต่อเนื่องไม่น้อยกว่า 3 ปี
\end{techbox}

\section{ความท้าทายเฉพาะทางในการจัดการสวนทุเรียนภาคตะวันออก}
\index{ความท้าทายเฉพาะทางในการจัดการสวนทุเรียนภาคตะวันออก}

การศึกษาและทำความเข้าใจเกี่ยวกับความท้าทายเฉพาะทางในการจัดการสวนทุเรียนภาคตะวันออก (Tropical Durian Orchard Challenges) มีความสำคัญอย่างยิ่งในระบบเกษตรอัจฉริยะ โดยมีรายละเอียดสาระสำคัญดังนี้

ทุเรียน (Durio zibethinus L.) เป็นพืชเศรษฐกิจที่มีความอ่อนไหวสูงต่อสภาพแวดล้อมอย่างยิ่ง โดยเฉพาะอย่างยิ่งปัญหาสำคัญในจังหวัดจันทบุรีและภาคตะวันออก:
- การเปลี่ยนแปลงสภาพภูมิอากาศอย่างฉับพลัน ฝนตกนอกฤดูกาลทำให้เกิดการแตกใบอ่อนแทนการสะสมตาดอก
- โรครากเน่าโคนเน่าจากเชื้อราไฟทอปธอรา (Phytophthora palmivora) เมื่อดินมีความชื้นแฉะสะสมเกินเกณฑ์
- ปัญหาน้ำต้นทุนขาดแคลนในช่วงพัฒนาการของผล (ระยะผลโตถึงผลแก่) ซึ่งหากขาดน้ำจะทำให้ทุเรียนสลัดผลทิ้งหรือเกิดอาการไส้ซึมเต่าเผา
- การตรวจวัดและจัดการด้วยเทคโนโลยี AIoT จึงเป็นทางออกเชิงโครงสร้างที่สามารถแก้ไขปัญหาเหล่านี้ได้อย่างตรงจุด

\section{ขั้นตอนและวิธีปฏิบัติการทดลอง}
\index{ขั้นตอนการทดลองประจำบทที่ 1}

\subsection{การสำรวจพื้นที่และวางแผนจุดติดตั้งโครงข่าย AIoT ในสวนทุเรียน}
\begin{sensorbox}[การสำรวจพื้นที่และวางแผนจุดติดตั้งโครงข่าย AIoT ในสวนทุเรียน]
1. สำรวจแผนผังแปลงสวนทุเรียน ระดับความสูงต่ำของพื้นที่ (Elevation) และตำแหน่งแหล่งน้ำ
2. ตรวจสอบระยะพิกัดสัญญาณและจุดอับสัญญาณ โดยใช้เครื่องมือทดสอบการครอบคลุมของคลื่นวิทยุ LoRaWAN (Field Tester)
3. กำหนดตำแหน่งติดตั้งสถานีตรวจอากาศหลัก (Weather Station) ในจุดที่เปิดโล่ง ไม่ถูกเงาต้นไม้บัง
4. วางจุดติดตั้งโพรบวัดดิน (Soil Probe Nodes) ใต้ทรงพุ่มของต้นทุเรียนต้นแบบตามระดับความสูงของดิน
5. ติดตั้งเกตเวย์รับส่งสัญญาณ (LoRaWAN Gateway) ณ ตำแหน่งอาคารศูนย์ปฏิบัติการที่มีเสาสูง
\end{sensorbox}

\subsection{การตรวจสอบการเชื่อมต่อข้อมูลระหว่างโหนดเซนเซอร์และคลาวด์เซิร์ฟเวอร์}
\begin{sensorbox}[การตรวจสอบการเชื่อมต่อข้อมูลระหว่างโหนดเซนเซอร์และคลาวด์เซิร์ฟเวอร์]
1. จ่ายไฟเลี้ยงให้โหนดเซนเซอร์ด้วยแผงโซลาร์เซลล์และแบตเตอรี่ LiFePO4
2. ตรวจสอบการส่งแพ็กเก็ตข้อมูลดิบ (Payload Uplink) ผ่านคลื่นความถี่ AS923 (มาตรฐานประเทศไทย)
3. เข้าสู่คอนโซลของ The Things Network (TTN) ตรวจดูค่า RSSI (Signal Strength) และ SNR (Signal-to-Noise Ratio)
4. เขียนสคริปต์ถอดรหัสเพย์โหลด (Payload Decoder) แปลงข้อมูลไบนารีเป็นค่าอุณหภูมิ ความชื้น และค่าแรงดันไฟฟ้า
5. ทดสอบส่งต่อข้อมูลไปยังแดชบอร์ดผ่านโพรโทคอล MQTT
\end{sensorbox}

\subsection{การตั้งค่าเงื่อนไขการแจ้งเตือนฉุกเฉินผ่าน Line Messaging API}
\begin{sensorbox}[การตั้งค่าเงื่อนไขการแจ้งเตือนฉุกเฉินผ่าน Line Messaging API]
1. สร้าง Line Notify Token หรือเชื่อมต่อผ่าน Line Official Account Developer Console
2. กำหนดเกณฑ์แจ้งเตือน: หากความชื้นสัมพัทธ์ในอากาศต่ำกว่า 40\% ร่วมกับอุณหภูมิสูงกว่า 38$^{\circ}$C ให้แจ้งเตือนภาวะเสี่ยงใบไหม้
3. ทดสอบจำลองค่าวิกฤตและตรวจสอบความเร็วในการส่งข้อความเตือนเข้าสู่กลุ่มผู้ดูแลสวน
\end{sensorbox}

\section{ตารางบันทึกผลการทดลองและการอภิปรายผล}
\begin{table}[H]
\centering\small
\caption{ตารางบันทึกผลการตรวจวัดและข้อมูลพารามิเตอร์ประจำบทที่ 1}
\label{tab:ch01_datasheet}
\begin{tabularx}{\textwidth}{p{0.3\textwidth} p{0.3\textwidth} X}
\toprule
\textbf{พารามิเตอร์ / การทดสอบ} & \textbf{ค่าที่ตรวจวัดได้} & \textbf{การแปลผลและการวิเคราะห์} \\
\midrule
จุดตรวจวัดที่ 1 (บริเวณดินดอน) & ค่าความชื้นอยู่ในเกณฑ์ปกติ & ปริมาณน้ำเพียงพอต่อการเจริญ \\
จุดตรวจวัดที่ 2 (บริเวณที่ลุ่ม) & มีการสะสมความชื้นสูง & ควรชะลอการให้น้ำเพื่อป้องกันรากเน่า \\
สถานะสัญญาณโครงข่าย & RSSI: -85 dBm, SNR: +8 dB & คุณภาพสัญญาณ LoRaWAN ยอดเยี่ยม \\
\bottomrule
\end{tabularx}
\end{table}

\section{คำถามทบทวนและแบบฝึกหัดท้ายบท}
\begin{enumerate}
    \item จงอธิบายความแตกต่างระหว่างระบบเกษตรอัจฉริยะ (Smart Agriculture) กับการทำเกษตรแบบดั้งเดิม
    \item เหตุใดเทคโนโลยี LoRaWAN จึงเหมาะสมกับการใช้งานในสวนไม้ผลพื้นที่กว้างมากกว่าสัญญาณ Wi-Fi ทั่วไป
    \item จงวิเคราะห์ความสัมพันธ์ระหว่างการแตกใบอ่อนของทุเรียนกับความผันผวนของความชื้นในดินและปริมาณฝน
\end{enumerate}
